\documentclass[11pt]{article}
\usepackage[margin=1in]{geometry}
\usepackage{amsmath,amssymb,amsthm}
\usepackage{graphicx}
\renewcommand{\labelitemi}{$\bullet$}
\renewcommand{\S}{Sec.~}
\newtheorem{theorem}{Theorem}
\newtheorem{proposition}[theorem]{Proposition}
\newtheorem{corollary}[theorem]{Corollary}
\newtheorem{lemma}[theorem]{Lemma}
\theoremstyle{remark}
\newtheorem{remark}[theorem]{Remark}
\newcommand{\Tr}{\operatorname{Tr}}
\title{The Three-Body Problem, Operator-First\\
\large A complete spectral anatomy of relational transport on the shape sphere,\\
a conserved record measured along the figure-eight,\\
and the shape manifold read as a birefringent crystal}
\author{Jeromie Beasley}
\date{August 22, 2026}
\begin{document}
\maketitle

\begin{abstract}
The planar three-body problem is treated as a question about one operator: the Gram matrix
of the differential of the map that sends a shape to its three mutual distances. Four
results follow. (i) That operator is coordinate-free and its singular values are
$\cos^2\vartheta$ and $1$ exactly, with $\vartheta$ the polar angle on the shape sphere and
no dependence on longitude -- replacing an earlier weighted tensor two of whose three
degeneracy loci were coordinate artefacts. (ii) The two singular loci are of different
species and were previously conflated: at syzygy the operator drops rank while its spectral
gap is \emph{maximal}, and at the equilateral configurations the gap closes while the
operator stays invertible. (iii) The equilateral degeneracy is \emph{not} a diabolical
point: the gap closes quadratically, $\mathrm{gap}/\vartheta^{2}\to1.000000$, the soft
direction returns with holonomy $+1$ rather than a sign, and the line field carries integer
index $+1$ at each pole -- so the obstruction to a global choice of most-rigid direction is
$\chi(S^{2})=2$ and there is no Berry phase to collect. (iv) Along the
Chenciner--Montgomery figure-eight the metric layer collapses at every one of six syzygies
per period, with $R=\sigma_{\min}^{-2}$ reaching $2.5\times10^{16}$, while the contour count
holds at $2$ to $5.1\times10^{-15}$ and the shifted trace-log equals
$\log(1+\sigma_{\min})+\log2$ to $2.2\times10^{-15}$ at every step. The unshifted ledger
fails, and fails for the stated reason: its branch point is enclosed by the contour. The
record outlives the metric on a real orbit, and the branch condition is doing visible work.
(v)~The open problem of the previous version is closed in the direction it predicted:
breaking the mass symmetry replaces the protected touching by a genuine conical
intersection, with the gap closing linearly and the eigenvector returning negated,
$-1.000000$ at every radius tested. (vi)~The two singular values are two index sheets, so
the degeneracies are optic axes, and the family is a crystal: equal masses are uniaxial
with one axis, unequal masses biaxial with exactly two, and a residual symmetry confines
both axes to its plane. A single point then admits a circle of ray directions rather than
one --- conical refraction, in shape space. (vii)~At a generic collinear configuration of
$N$ bodies the rank of $d\Phi$ is $N-2$ for a reason available in one line, so the
three-body rank drop is the smallest member of a family.
\end{abstract}

\section{What is and is not being solved}

No closed-form solution is offered and none is possible. What is offered is a complete and
coordinate-free spectral anatomy of \emph{relational transport} -- how much the three mutual
distances resist change per unit motion of the shape -- together with two quantities that
survive the places where that anatomy degenerates. The claim to novelty is narrow: the
singular behaviour of side-length coordinates at collinear configurations is ninety years
old (Murnaghan 1936; Lema\^itre 1952), and what is new here is its statement as a rank drop
of an explicit operator, the correct separation of two loci that the literature and my own
earlier draft treated as one, and the measurement of a conserved record through them on an
actual orbit.

\section{The operator, and why this one}

Let $S^2$ be the Kendall shape sphere of the planar three-body problem in the coordinates
$w=(\sin\vartheta\cos\varphi,\sin\vartheta\sin\varphi,\cos\vartheta)$, with $w_3$ the signed
area, so the equator $w_3=0$ is the collinear (syzygy) locus and the poles are the
equilateral (Lagrange) configurations.

The three squared mutual distances $r_{ij}^2$ are affine in $(w_1,w_2)$ and independent of
$w_3$. Write $\Phi:S^2\to\mathbb R^2$ for that map. The object of study is
\begin{equation}
M \;=\; d\Phi^{\!\top}d\Phi ,
\label{eq:M}
\end{equation}
the Gram matrix of its differential, read in an orthonormal frame of the round metric.

\begin{theorem}[Exact spectrum]\label{t:spec}
In the frame $\{e_\vartheta,e_\varphi\}$,
$d\Phi(e_\vartheta)=(\cos\vartheta\cos\varphi,\cos\vartheta\sin\varphi)$ and
$d\Phi(e_\varphi)=(-\sin\varphi,\cos\varphi)$, so
\[
M=\begin{pmatrix}\cos^{2}\vartheta&0\\[2pt]0&1\end{pmatrix},
\qquad
\sigma_{\min}=\cos^{2}\vartheta,\quad\sigma_{\max}=1 ,
\]
with no dependence on $\varphi$ and no free constant.
\end{theorem}
\begin{proof}
Direct substitution of the embedding derivatives into \eqref{eq:M}; the cross term is
$-\cos\vartheta\cos\varphi\sin\varphi+\cos\vartheta\sin\varphi\cos\varphi=0$.
\end{proof}

\begin{remark}[What this replaces, and why the replacement was forced]
An earlier draft used the weighted tensor $T(w)=\operatorname{diag}(w_1^2,w_2^2,w_3^2)$.
That object is \emph{axis-dependent}: its determinant on the sphere is
$\tfrac34\sin^6\vartheta\cos^2\vartheta\sin^2 2\varphi$, which vanishes on the great circles
$\varphi=0$ and $\varphi=\pi/2$ at generic latitude -- two degeneracy loci with no invariant
meaning, since only $w_3=0$ is coordinate-free. Theorem~\ref{t:spec} keeps the one locus
that survives a change of axes and discards the two that do not. The relational rigidity
\[
R=\sigma_{\min}^{-2}=\sec^{4}\vartheta
\]
is unchanged by the replacement, and now has prefactor exactly $1$ and no longitude
dependence, where the earlier route reported a $\varphi$-averaged constant.
\end{remark}

\section{Two degeneracies, of different species}

\begin{proposition}[Rank drop at syzygy; the gap is maximal there]\label{p:syzygy}
At $\vartheta=\pi/2$ the smaller eigenvalue of $M$ vanishes while the larger equals $1$, so
the spectral gap is $1-\cos^{2}\vartheta=\sin^{2}\vartheta$, which attains its
\emph{maximum} exactly on the syzygy locus.
\end{proposition}

\begin{proposition}[Degeneracy at the equilateral points; the operator stays invertible]
At $\vartheta\to0,\pi$ both eigenvalues tend to $1$; the gap closes while $\det M\to1$.
\end{proposition}

\noindent
These are opposite situations and the distinction is not cosmetic. A Riesz projector
requires a gap, not invertibility: it continues real-analytically through syzygy, where the
gap is widest, and it is only at the equilateral points that the two spectral subspaces
cease to be separable. Any ``diabolical point'' or projector-collapse language belongs at
the poles, and any regularisation introduced to carry a projector across the equator is
unnecessary -- there is nothing there to regularise.

\section{The equilateral degeneracy is not diabolical}

A generic degeneracy of a real symmetric $2\times2$ family in two parameters is a conical
intersection: the gap closes linearly, the eigenvector line field winds by half a turn, and
one circuit returns the eigenvector multiplied by $-1$. None of that happens here.

\begin{theorem}[Quadratic closing, integer index, no Berry phase]\label{t:lagrange}
Near a pole the gap closes quadratically,
$\mathrm{gap}/\vartheta^{2}\to1$; the softest eigendirection transported once around the
pole returns with holonomy $+1$; and as a line field on the sphere it has index $+1$ at each
pole.
\end{theorem}
\begin{proof}[Proof and verification]
$\mathrm{gap}=1-\cos^{2}\vartheta=\sin^{2}\vartheta$, so the ratio to $\vartheta^{2}$ tends
to $1$: measured $0.996671$, $0.999967$, $1.000000$, $1.000000$ at
$\vartheta=10^{-1},10^{-2},10^{-3},10^{-4}$. The soft direction is $e_\vartheta$, which is
smooth and single-valued along any circle of constant latitude, so the transported vector
returns to itself: measured holonomy $+1.000000$ at $\vartheta=0.3,0.05,0.01$ with zero sign
flips during continuation. Its winding in the coordinate frame is $0$ and, as a field on the
sphere, $+1$ -- the coordinate frame supplies the whole rotation
(\texttt{LAGRANGE-PHASE-01}).
\end{proof}

\begin{corollary}[The obstruction is topological, not local]
Two poles of index $+1$ sum to $2=\chi(S^{2})$. There is therefore no global smooth choice
of ``most rigid direction'' on the shape sphere, and the obstruction is the Euler
characteristic rather than any stratum.
\end{corollary}

\begin{remark}[Why the linear term is absent]
The equilateral configuration is fixed by the $Z_3$ relabelling of the three bodies, and an
invariant symmetric $2$-tensor at a fixed point of that action must be a multiple of the
identity to leading order. The linear anisotropy is forbidden by symmetry, which is why the
degeneracy is quadratic and the winding integer rather than half-integer. A generic
three-body problem with unequal masses breaks $Z_3$ and is the place to look for a genuine
conical intersection.
\end{remark}

\section{The $\cos^{4}$ law is a fold, not a three-body fact}

\begin{proposition}[$A_2$ normal form]\label{p:fold}
Near the syzygy locus write $\vartheta=\pi/2+\epsilon$. Then
$w_1=\cos\epsilon\cos\varphi\sim(1-\epsilon^{2}/2)\cos\varphi$: the displacement of the
image of $\Phi$ is quadratic in $\epsilon$, which is Whitney's fold $x\mapsto x^{2}/2$. The
singular value of $d\Phi$ has a simple zero, the Gram matrix squares it, and $R$ squares it
again, giving $\sec^{4}$.
\end{proposition}

\noindent
So the exponent $4$ is two bookkeeping squarings on top of a simple zero, and the invariant
content of the ``$\cos^4$ rigidity law'' is only that $d\Phi$ has a simple zero on a
codimension-one locus. Since a fold is the generic singularity of a map between
equal-dimensional manifolds, the same exponent must appear for any relational rigidity
defined this way at a generic degenerate locus. That converts a recurring coincidence into a
theorem and predicts where it should fail: at an $A_3$ cusp, which the equal-mass problem
does not have and which unequal masses or four bodies may supply.

\section{The transport form is pinched, and by two independent mechanisms}

The shape sphere is the $d=2$ rung of the square case $N=d+1$, whose transport operator on
the eigenvalue block is $A=D_t-g\,t^{\!\top}$.

\begin{theorem}[Reciprocity]\label{t:recip}
$A$ is complex-symmetric with respect to the bilinear pairing
$B=\operatorname{diag}(t/g)$: $A^{\!\top}B=BA$ exactly, with left eigenvectors $\ell_n=BR_n$.
Consequently, for probes symmetric under the same pairing the outbound and return
amplitudes coincide, the round-trip form is a perfect square, the rank measure $W$ vanishes,
and $\det g\le0$ with equality when the amplitudes are phase-aligned.
\end{theorem}

\noindent
This is the symmetrisation theorem $D_t-g\,t^{\!\top}=S(D_t-uu^{\!\top})S^{-1}$ read as a
pairing rather than a similarity, with $B=S^{-2}$. Because the operator, the pairing and the
spectrum are all real on the regular stratum, the amplitudes are real, their relative phase
is zero, and $\det g=0$ \emph{exactly}: measured $W=1.7\times10^{-16}$ and $\det g=0.00$ with
$B$-symmetric probes, against $W=0.17$ and $\det g=-1.5\times10^{-4}$ with generic probes on
the same operator (\texttt{BODY-RECIP-01}). The transport form carries no cone anywhere on
the regular stratum, and $\|B\|=\max_i t_i/g_i$ diverges exactly on the collision divisor --
so the only place the geometry could open is the boundary.

\section{The record along the figure-eight}

\begin{figure}[t]\centering
\includegraphics[width=0.94\textwidth]{fig_fig8.pdf}
\caption{The two layers along one period of the Chenciner--Montgomery figure-eight
(\texttt{FIG8-LEDGER-01}). Above: the signed area $w_3$, whose six zeros are the syzygies.
Middle: the metric layer $R=\sigma_{\min}^{-2}$, spiking through sixteen orders at every
crossing. Below: the contour count and the shifted trace-log, flat through all six.}
\label{f:fig8}
\end{figure}

Everything above is a statement about the shape sphere. The orbit supplies the test: the
figure-eight walks the system through the degenerate locus repeatedly and unassisted.

\paragraph{Protocol.} Equal masses, the standard initial data, integrated by fourth-order
Runge--Kutta at $h=2\times10^{-5}$ over one period $T=6.32591398$; closure
$|q(T)-q(0)|=4.1\times10^{-5}$. At each step the configuration is mapped to the shape sphere
and the spectrum $\{\cos^{2}\vartheta,1\}$ of Theorem~\ref{t:spec} is formed. Two contour
quantities are then evaluated on a circle of radius $0.9$ about $0.5$, which encloses both
eigenvalues with margin: the count $N=\frac{1}{2\pi i}\oint\Tr R(z)\,dz$ and the shifted
trace-log $Z=\frac{1}{2\pi i}\oint\log(z+\eta)\Tr R(z)\,dz$ with $\eta=1$.

\begin{theorem}[Measured: the record outlives the metric]\label{t:fig8}
Along one period the orbit crosses syzygy six times -- confirmed independently by the sign
changes of the raw triangle area -- and
\begin{itemize}
\item the metric layer collapses at every crossing: $\sigma_{\min}$ falls to
$6.3\times10^{-9}$ at the sampled points and to $0$ at the crossings themselves, with
$R=\sigma_{\min}^{-2}$ reaching $2.5\times10^{16}$;
\item the count holds: $\max|N-2|=5.1\times10^{-15}$ over the whole orbit, crossings
included;
\item the shifted trace-log holds: $Z=\log(1+\sigma_{\min})+\log2$ to
$2.2\times10^{-15}$ at every step, finite everywhere, taking its minimum $\log2$ exactly at
each syzygy.
\end{itemize}
\end{theorem}

\begin{remark}[The branch condition is not a technicality]
Repeating the computation with $\eta=0$ breaks it, and breaks it for the stated reason: the
branch point of the logarithm then sits at the origin, which is \emph{inside} the contour,
so the integrand is not meromorphic in the enclosed region and the result is not the sum of
the enclosed logarithms (measured minimum $-0.580$, against $\log2=0.693$ for the correct
evaluation). The hypothesis under which the ledger is defined -- a branch analytic on and
inside $\Gamma$ -- is doing visible work here, and the orbit shows what its failure looks
like. It is a branch condition, not a positivity condition, and the shift is what restores
it.
\end{remark}


\section{The symmetry broken: a genuine conical intersection}

Theorem~\ref{t:lagrange} attributes the trivial holonomy to $Z_3$ and therefore predicts
that breaking the mass symmetry should produce the generic alternative. It does.

\begin{theorem}[Conical intersection at unequal masses]\label{t:conical}
For $m=(1,2,3)$ the two singular values coincide at isolated points of the shape sphere.
At one of them, $a=0.845154255$, $\psi=5.034139535$, both equal $\sqrt2$ and the gap closes
to $1.5\times10^{-14}$. On a circle of radius $r$ about that point the gap is
\emph{linear}, $\mathrm{gap}/r = 2.16437,\,2.16294,\,2.16281,\,2.16280$ at
$r=10^{-2},3\times10^{-3},10^{-3},3\times10^{-4}$, and the softest eigenvector transported
once around returns \emph{negated}: holonomy $-1.000000$ at every radius tested.
\end{theorem}

\noindent
Linear closing, sign flip, Berry phase $\pi$: the generic real-symmetric degeneracy in two
parameters, appearing exactly when the protection is removed. The prediction was recorded
before the test (\texttt{UNEQUAL-CHECK-01}).

\section{The shape manifold is a birefringent crystal}

The two singular values of $M$ are two sheets of an index surface, so a point where they
coincide is an \emph{optic axis}, and the classification of such points is the
classification of crystals. That correspondence is exact here, not decorative.

\begin{theorem}[Uniaxial, biaxial, and the symmetry plane]\label{t:crystal}
Counting coincidences over the whole sphere (\texttt{BIAXIAL-01}):
equal masses give two points, $a=1/\sqrt2$ at $\psi=\pi/2,3\pi/2$, i.e.\ \emph{one} axis
meeting the sphere twice --- the uniaxial case, with the touching forced by symmetry.
Masses $(1,2,3)$ give \emph{four} points in two pairs, at $a^2=5/7$ and $a^2=2/7$, both
carrying $\sigma_1=\sigma_2=\sqrt2$ --- two axes, the biaxial case. Masses $(1,1,5/2)$,
which retain one reflection, also give four points, but every one lies at
$\psi=\pi/2$ or $3\pi/2$: the axes are confined to the symmetry meridian.
\end{theorem}

\begin{remark}[The crystallographic reading, and what it predicts]
A uniaxial crystal has its sheets touching along a single symmetry-forced axis; a biaxial
crystal has no such symmetry and exactly two isolated optic axes; and a crystal retaining
one symmetry plane has both axes lying in it. The three rows of Theorem~\ref{t:crystal} are
those three statements. The prediction that transfers is Hamilton's: at a biaxial optic
axis the sheets meet conically, so the surface normal --- the ray direction --- is
undefined there, and a single incident ray refracts into a hollow cone.
\end{remark}

\begin{proposition}[A circle of ray directions]\label{p:ray}
Near an optic axis write the traceless part of $M$ as a linear map $J$ of the two shape
coordinates; $\det J=-4.463$, so the cone is nondegenerate. Approaching the axis from
twelve equally spaced directions, the gradient of the lower sheet points in ten distinct
directions, with $|\nabla\lambda|$ ranging over $0.638$ to $8.080$
(\texttt{CONICAL-RAY-01}).
\end{proposition}

\noindent
So a single point of shape space admits a circle of ray directions rather than one. This is
the only statement in the paper about what a \emph{trajectory} would do rather than what an
operator's spectrum is, and it is the natural next computation.

\begin{figure}[t]\centering
\includegraphics[width=0.96\textwidth]{fig_crystal.pdf}
\caption{The crystal reading. Left and centre: coincidences of the two singular values on
the shape sphere, two points for equal masses (one optic axis) and four for $m=(1,2,3)$
(two axes). Right: the ray direction against the direction of approach at one axis --- the
gradient does not converge, which is what makes the refraction conical.}
\label{f:crystal}
\end{figure}

\section{The rank drop is the smallest member of a family}

\begin{proposition}[Collinear rank for $N$ bodies]\label{p:rankN}
At a generic collinear configuration of $N$ planar bodies, $\operatorname{rank}d\Phi=N-2$,
so the kernel also has dimension $N-2$ in the $(2N-4)$-dimensional shape space.
\end{proposition}
\begin{proof}
Squared distances depend on the transverse offsets quadratically:
$\partial r_{ij}^2/\partial y_k \propto (y_i-y_j)$, which vanishes identically when all
$y=0$. Every transverse direction is therefore in the kernel, and the transverse
displacements number $N-2$ after removing translation and the rotation that acts as
$y_i\propto x_i$.
\end{proof}

\noindent
\emph{Verified} for $N=3,\dots,8$: ranks $1,2,3,4,5,6$ against shape dimensions
$2,4,6,8,10,12$ (\texttt{LADDER-CHECK-01}). At $N=3$ a single direction collapses, which is
the syzygy result of \S3; for larger $N$ several collapse at once, so the three-body case is
the unusually clean member rather than the representative one. The regular polygons carry
structured degeneracies of their own --- the square gives $\{\sqrt8,2,2,2\}$ with a
threefold degeneracy that splits linearly under perturbation, the pentagon
$\{\sqrt{5+\sqrt5},\sqrt{5+\sqrt5},2,2,\sqrt{5-\sqrt5},\sqrt{5-\sqrt5}\}$ in three pairs.

\section{Transient amplification of the figure-eight}\label{s:transient}

The orbit is linearly stable, which is a statement about eigenvalues. Non-normal operator
theory supplies a second, independent rate, and it is the one that governs finite
perturbations over one period.

\begin{proposition}[Measured]\label{p:transient}
In mass-weighted Jacobi coordinates --- where translation and boost are absent by
construction and the Euclidean norm is the kinetic-energy norm --- perturbations restricted
to the energy and angular-momentum level set and measured modulo the time-shift and
rotation directions are amplified by $45.7$ in energy, $6.8\times$ in amplitude, within one
period, returning to $9.6$ at the end. Energy is conserved to $6.5\times10^{-14}$ and
angular momentum to $2.2\times10^{-14}$ (\texttt{TB-TRANSIENT-03}).
\end{proposition}

\noindent
The amplification is timed by the shape: the mean gain in the quartile of the orbit closest
to syzygy is $10.58$, and in the quartile farthest from it $26.10$. The burst lives on the
triangular configurations, not at the wall --- which is what \S6 requires, since at syzygy
the transport form is degenerate and there is no geometry left to amplify through.

\begin{remark}[What this is not]
The quotient used is an orthogonal complement of the symmetry and integral directions, and
that complement is not invariant under the flow. The object measured is therefore a
compression of the propagator, not a restriction, and its eigenvalues are \emph{not} the
Floquet multipliers; only the gain is quoted. The practical reading is nonetheless
available: a perturbation of size $\epsilon$ reaches $\approx7\epsilon$ before the neutral
multipliers return it, so the orbit's usable basin is set by that excursion rather than by
its linear stability.
\end{remark}

\section{What is claimed}

Not a solution of the three-body problem. Not a new integral of motion in the classical
sense: $N$ and $Z$ are functions of the shape alone, and $Z$ varies smoothly along the orbit
between $\log 2$ and $1.045$. What is claimed is that they are \emph{insensitive to the
degeneracy that destroys the metric layer}, at machine precision, on a real solution, six
times per period -- and that this is the same ordering measured on synthetic operator
families, now realised in celestial mechanics with no dial to turn by hand.

The remaining content is the anatomy: an invariant operator with an exact spectrum, two
degeneracies correctly separated, an equilateral degeneracy that is quadratic rather than
conical and therefore carries no geometric phase, a $\cos^4$ law identified as a fold normal
form and therefore universal, a transport form pinched everywhere on the regular stratum by
two independent mechanisms, a genuine conical intersection the moment the mass symmetry is
broken, the uniaxial/biaxial classification with its optic axes located, the collinear rank
law for all $N$, and a transient amplification of the figure-eight that its linear stability
does not predict.

\section{Open}

\begin{enumerate}
\item The relation $a_1^2+a_2^2=1$ between the two optic axes, and the exact rationals
$a^2=5/7,\,2/7$ at $m=(1,2,3)$ against a total mass of $6$. Observed to the printed digits;
not derived.
\item A closed form for the regular-$N$-gon spectra. The hexagon's spectrum contains the
pentagon's values, which a table of special cases would not do.
\item Proposition~\ref{p:fold} predicts a different exponent at an $A_3$ cusp. The
equal-mass shape sphere has folds only; where the first cusp appears is not known.
\item The ledger \emph{on} the divisor, where the pairing $B$ of Theorem~\ref{t:recip} is
undefined and the count is the only surviving object.
\item A genuinely reduced monodromy for \S\ref{s:transient}: the quotient used there is an
orthogonal complement and therefore a compression, not a restriction, so its spectrum is
not the Floquet spectrum. The symplectic complement inside $\ker dE\cap\ker dL$ would make
it one.
\item Conical refraction at the level of trajectories rather than of ray directions:
integrate an orbit into an optic axis and ask what the reduced dynamics does there.
\end{enumerate}

\end{document}
